\chapter{Architectural Design}
% Overall system architecture
\section{System Overview}
\begin{figure}[ht]
\centering
\resizebox{\linewidth}{!}{%
\begin{tikzpicture}[
    font=\sffamily,
    >=Latex,
    node distance=14mm and 16mm,
    proc/.style={
        draw,
        rounded corners=2pt,
        rectangle,
        minimum height=9mm,
        minimum width=31mm,
        align=center,
        line width=0.9pt,
        fill=blue!6,
        draw=blue!55!black
    },
    data/.style={
        draw,
        rounded corners=2pt,
        rectangle,
        minimum height=9mm,
        minimum width=31mm,
        align=center,
        line width=0.9pt,
        fill=green!7,
        draw=green!45!black
    },
    outbox/.style={
        draw,
        rounded corners=2pt,
        rectangle,
        minimum height=9mm,
        minimum width=31mm,
        align=center,
        line width=0.9pt,
        fill=orange!10,
        draw=orange!60!black
    },
    eval/.style={
        draw,
        rounded corners=2pt,
        rectangle,
        minimum height=10mm,
        minimum width=36mm,
        align=center,
        line width=1.0pt,
        fill=purple!8,
        draw=purple!55!black
    },
    line/.style={-Latex, line width=0.9pt, draw=black!75},
]

% Left inputs and mixture generator
\node[data] (A) {Speaker Audio\& Transcripts};
\node[data, below=18mm of A] (C) {Noise Audio};
\node[proc, right=28mm of $(A)!0.5!(C)$] (B) {Mixture Generator};

% Direct outputs from mixture generator
\node[outbox, right=28mm of B] (D) {Mixed Audio};
\node[data, below=14mm of B] (J) {Manifest};

% ASR models
\node[proc, right=28mm of D, yshift=14mm] (E) {ASR Model 1\\(e.g. Whisper)};
\node[proc, right=28mm of D] (F) {ASR Model 2\\(Wav2Vec2)};
\node[proc, right=28mm of D, yshift=-14mm] (G) {ASR Model 3\\(Parakeet)};

% Scoring and final results
\node[eval, right=30mm of F] (H) {Scoring:\\WER, cpWER};
\node[outbox, below=18mm of H] (I) {Results and Analysis};

% Edges from inputs to generator
\draw[line] (A) -- (B);
\draw[line] (C) -- (B);

% Edges from generator
\draw[line] (B) -- (D);
\draw[line] (B) -- (J);

% Mixed audio to ASR models
\draw[line] (D) -- (E);
\draw[line] (D) -- (F);
\draw[line] (D) -- (G);

% ASR outputs and manifest to scoring
\draw[line] (E) -- (H);
\draw[line] (F) -- (H);
\draw[line] (G) -- (H);
\draw[line] (J) -| (H.west);

% Final output
\draw[line] (H) -- (I);

\end{tikzpicture}%
}
\caption{System data flow: input speech and noise files (.wav or .flac) are combined into mixtures, which are then transcribed by ASR models. The transcripts are scored against the reference with WER and overlap-aware metrics.}
\label{fig:general_flow}
\end{figure}

\Cref{fig:general_flow} depicts the high-level data flow through the system. Inputs include clean speaker audio folders (containing audio of a single speaker, with speaker ID as the folder name and the corresponding transcripts) and noise audio folders (.wav or .flac, with noise type as the folder name). A YAML configuration file specifies parameters for mixture generation, including the number of mixtures, SNR levels (SNR), overlap ratios (ovr), duration, and speaker count (number of unique speakers). The \emph{Mixture Generator} module creates synthetic mixtures based on these configuration and outputs both the audio files and a manifest (containing metadata and transcripts). The mixed audio is fed into the target ASR models for inference. The result hypotheses are scored against the reference transcripts using evaluation metrics. The resulting hypotheses are then evaluated against the reference transcripts using specified evaluation metrics. The final results is analysed and compiled into a final report.

% Synthetic mixture generation process
\section{Synthetic Mixture Generation}
The mixture generation process comprises three steps: (1) \textbf{Speaker and clip selection}: randomly select the required number of speakers and clips from the clean speech dataset, (2) \textbf{Offset scheduling}:
(1) \textbf{Offset computation}: calculate the time offsets for each utterance such that the overall overlap ratio approaches the target. (2) \textbf {Time offset adjustment}: adjust the time offsets to ensure the overall overlap ratio approaches the target. (3) \textbf{Noise addition}: add noise to the mixture at the specified SNR level. The output of this process will be a set of audio with different speaker counts and overlap ratios as base clips, and each base clip will have variants with different SNR levels. The metadata and transcripts of the mixtures are stored in a manifest file.

\subsection{Speaker and Clip Selection}

\begin{figure}[ht]
\centering
\resizebox{\linewidth}{!}{%
\begin{tikzpicture}[
    font=\sffamily,
    >=Latex,
    node distance=10mm and 16mm,
    every node/.style={align=center},
    process/.style={
        rectangle,
        rounded corners=2pt,
        draw=blue!60!black,
        very thick,
        minimum width=6.2cm,
        minimum height=10mm,
        text width=6.0cm,
        fill=blue!5
    },
    decision/.style={
        diamond,
        aspect=2.2,
        draw=orange!70!black,
        very thick,
        minimum width=5.4cm,
        minimum height=11mm,
        text width=3.4cm,
        fill=orange!10,
        inner sep=1pt
    },
    line/.style={
        -{Latex[length=3mm,width=2mm]},
        very thick,
        draw=black!75
    },
    labelstyle/.style={font=\sffamily\footnotesize, inner sep=1pt, fill=white, text=black!80}
]

% Nodes
\node[process, text width=4.4cm, minimum width=4.8cm] (A) {Speaker count ($c$) \& duration ($d$) \\ \& overlap ratio ($o$)};
\node[process, right=of A, text width=4.0cm, minimum width=4.5cm] (B) {randomly choose ($c$) speakers};
\node[process, right=of B, text width=4.0cm, minimum width=4.5cm] (C) {randomly choose one speaker};
\node[process, right=of C, text width=4.2cm, minimum width=4.6cm] (D) {randomly choose one clip from the speaker};
\node[decision, right=of D, text width=3.8cm, minimum width=4.3cm] (E) {if total clip duration $> d\,(1+o)$};
\node[process, right=of E, text width=5.4cm, minimum width=5.8cm] (F) {return sequence of clips + sequence of transcription};

% Main flow
\draw[line] (A) -- (B);
\draw[line] (B) -- (C);
\draw[line] (C) -- (D);
\draw[line] (D) -- (E);
\draw[line] (E) -- node[labelstyle, right] {True} (F);

% False loop back
\draw[line] (E.south) -- ++(0,-8mm) -| node[labelstyle, pos=0.2, above] {False} (C.south);

\end{tikzpicture}%
}
\caption{Speaker and Clip Selection Process}
\label{fig:clip_selection}
\end{figure}
This process takes input of the target speaker count $c$ and duration constraints $d$ (minimum duration of final mixed audio), and returns: (1) a sequence of clips satisfying the target speaker count and duration constraints, (2) the corresponding sequence of transcripts.

\Cref{fig:clip_selection} shows the process that is designed to create a sequence of clips with specified speaker count and duration. A set of $c$ speakers is selected at uniform random without replacement. During the clip selection process, the next speaker is chosen with probability inversely proportional to the number of clips contributed from each speaker to balance the number of clips per speaker (with probability of 0 for speakers from the previous iteration). The random selection of clips is uniform random with replacement. The process iterates until the total duration of the selected clips exceeds the target duration adjusted by the overlap ratio (to account for the overlapping region that is double-counted here). The adjusted total clip duration $T_{sum}$ is calculated by:
\[
    \begin{aligned}
    \mathrm{OVR} &= T_{overlap} / d \\
    T_{duration} &= T_{sum} - T_{overlap} \\
    T_{overlap} &= T_{duration} - d \\
    \mathrm{OVR} &= (T_{sum} - d) / d \\
    \Rightarrow\quad T_{sum} &= d \cdot (1 + \mathrm{OVR})
    \end{aligned}
\]
The output transcript is a sequence of speaker-attributed transcripts.

\subsection{Offset Scheduling}
\begin{figure}[ht]
\centering
\begin{minipage}[t]{0.48\linewidth}
\vspace{0pt}
\centering
\resizebox{\linewidth}{!}{
\begin{tikzpicture}[
    font=\sffamily,
    >=Latex,
    node distance=14mm and 14mm,
    line/.style={draw, -{Latex[length=3mm,width=2mm]}, very thick, color=gray!70},
    block/.style={
        rectangle,
        rounded corners=2pt,
        draw=blue!50!black,
        very thick,
        fill=blue!5,
        text width=36mm,
        minimum height=11mm,
        align=center,
        inner sep=6pt
    },
    decision/.style={
        diamond,
        aspect=2,
        draw=orange!70!black,
        very thick,
        fill=orange!10,
        text width=20mm,
        align=center,
        inner sep=2pt
    },
    labelstyle/.style={font=\footnotesize\bfseries, text=gray!35!black, fill=white, inner sep=1pt}
]

% Nodes
\node[block] (A) {sequence of clips\\[-1pt] \& overlap ratio target $(\mathrm{ovr})$\\[-1pt] \& clip spacing mean $(\mu)$\\[-1pt] \& clip spacing std $(\mathrm{std})$};
\node[decision, below=of A] (B) {if $(\mathrm{ovr}) > 0$};
\node[block, below left=12mm and 11mm of B, text width=28mm] (C) {overlap allocation (\Cref{fig:overlap_allocation})};
\node[block, below right=12mm and 11mm of B, text width=31mm] (D) {allocate spacing\\[-1pt] $=\,\mathrm{Norm}(\mu,\mathrm{std})$};
\node[block, below=16mm of $(C)!0.5!(D)$, text width=34mm] (E) {sequence of clip start times};

% Edges
\draw[line] (A) -- (B);
\draw[line] (B) -- node[labelstyle, left, pos=0.45] {True} (C);
\draw[line] (B) -- node[labelstyle, right, pos=0.45] {False} (D);
\draw[line] (C) -- (E);
\draw[line] (D) -- (E);

\end{tikzpicture}
}
\end{minipage}\hfill
\begin{minipage}[t]{0.48\linewidth}
\vspace{0pt}
\centering
\resizebox{\linewidth}{!}{
\begin{tikzpicture}[
    font=\sffamily,
    >=Latex,
    node distance=10mm and 16mm,
    line/.style={draw=gray!70, -{Latex[length=3mm,width=2mm]}, thick},
    box/.style={
        rectangle,
        rounded corners=2pt,
        draw=blue!55!black,
        fill=blue!6,
        thick,
        align=center,
        text width=8.2cm,
        minimum height=10mm,
        inner sep=4pt
    },
    decision/.style={
        diamond,
        aspect=2.1,
        draw=orange!70!black,
        fill=orange!10,
        thick,
        align=center,
        text width=3.2cm,
        inner sep=1.5pt
    }
]

% Nodes
\node[box] (A) {sequence of clips \& overlap ratio $(\mathrm{ovr})$ \& overlap std ratio};

\node[box, below=of A] (B) {abs size of overlap $(o_{\mathrm{abs}})$\\
$= (\mathrm{ovr}) \cdot \dfrac{\text{total length of all clips}}{1 + (\mathrm{ovr})}$};

\node[box, below=of B] (C) {average abs overlap $(\bar{o}_{\mathrm{abs}})$\\
$= \dfrac{o_{\mathrm{abs}}}{\text{number of clips} - 1}$};

\node[box, below=of C] (D) {sequence of overlap\\
$= \mathrm{Norm}(\bar{o}_{\mathrm{abs}},\, \bar{o}_{\mathrm{abs}} \cdot \mathrm{std\_ratio})$};

\node[box, below=of D] (E) {difference $(d)$\\
$= \left|\left(\sum \text{sequence of overlap}\right) - o_{\mathrm{abs}}\right|$};

\node[decision, below=12mm of E] (F) {if $d=0$};

\node[box, below left=14mm and 28mm of F, text width=4.0cm, fill=green!8, draw=green!50!black] (G) {return sequence of overlap};

\node[box, below right=14mm and 28mm of F, text width=4.6cm, fill=red!6, draw=red!55!black] (H) {allocate $(d)$ at uniform random};

% Arrows
\draw[line] (A) -- (B);
\draw[line] (B) -- (C);
\draw[line] (C) -- (D);
\draw[line] (D) -- (E);
\draw[line] (E) -- (F);

\draw[line] (F) -- node[midway, left, font=\footnotesize\sffamily] {True} (G);
\draw[line] (F) -- node[midway, right, font=\footnotesize\sffamily] {False} (H);
\draw[line] (H) -- (G);

\end{tikzpicture}
}
\caption{Overlap allocation process when overlap ratio $> 0$}
\label{fig:overlap_allocation}
\end{minipage}
\caption{Offset Scheduling Process}
\label{fig:offset_scheduling}
\end{figure}
This process takes input of a sequence of clips, overlap ratio ($\mathrm{ovr} \in [0,1]$), overlap std ratio, and clip spacing mean and std (for non-overlapping case) and outputs the start sample index for each clip where the overlap ratio is exactly as specified.

In the case where the overlap ratio is zero, the start time of each clips is determined by allocating a spacing that is a random sample from a normal distribution \texttt{Norm}($\mu$, $\mathrm{std}$) where $\mu$ and $\mathrm{std}$ are the mean and std of the clip spacing specified.
In the case where overlap ratio is greater than zero, we first compute the target overlap $o_{\mathrm{abs}}$ in samples by:
\[
o_{\mathrm{abs}} = \frac{\mathrm{ovr}}{1+\mathrm{ovr}} \cdot \sum \mathrm{clip~duration}.
\]
and, for $n$ number of clips, the average overlap $\bar{o}_{\mathrm{abs}}$ by $\bar{o}_{\mathrm{abs}} = \frac{o_{\mathrm{abs}}}{n-1}$. Then each overlap between adjacent clips is then sampled from a normal distribution \texttt{Norm}($\bar{o}_{\mathrm{abs}}$, $\bar{o}_{\mathrm{abs}} \cdot \mathrm{std\_ratio}$) where the overlap is capped by the duration of the preceding and succeeding clips. After that, a residual correction is applied so that the sum of sampled overlaps equals the target overlap exactly. The difference $\mathrm{diff} = o_{\mathrm{abs}} - \sum\text{sampled overlaps}$. If the difference is greater than zero, we randomly select overlaps that are smaller than the preceding clip and succeeding clip duration, and add one sample to the selected overlaps until the difference is zero. If the difference is smaller than zero, we randomly select overlaps that are greater than zero and subtract one sample until the difference is zero.
Finally, the start sample index of each clip is determined by adding the spacing (overlap will be 
treated as negative spacing). The start and end timestamp of each clip is then put into the manifest.
\subsection{Noise Addition}
The noise addition process takes input of the mixed audio, set of noise audio, noise type and target SNR level ($\mathrm{SNR}_{dB} \in \mathbb{R} ~\text{or}~ \text{clean}$, measured over active speech time), and output the noisy mixed audio, where same base audio can produce different noisy versions. The noise type can be domestic (D), public (P), transportation (T), nature (N), office (O), street (S) or all (A) (more details in \Cref{sec:dataset}). Then sample noise audio from the specified noise type(s) uniformly at random with replacement, and concatenate the noise audio until the total length is reached (exceeding part will be truncated). Finally, the noise is scaled by a factor $\alpha$ to meet the target SNR level. The scaling factor $\alpha$ is computed by:
\[
\alpha = \frac{A_\text{active speech}}{A_\text{noise} \times 10^{\mathrm{SNR}_{dB}/20}}.
\]
where $A_\text{active speech}$ is the root mean square (RMS) amplitude of the active speech region in the mixture, and $A_\text{noise}$ is the RMS amplitude of the noise segment before scaling. The active speech region is determined by the RMS of all clips in the mixture. The SNR is measured relative to the active speech region to ensure the noise level is accurately proportional to the speech content.
Finally, the scaled noise is added to the mixture to produce the final audio.

This algorithm is designed to support controlled experiments on overlap ratio, SNR level, speaker count, duration and noise type by varying each factor independently. It also supports noise analysis on the same clip with different noise types and SNR levels.

% metadata and transcription structure
\subsection{Manifest Structure}
The manifest is a csv file that contains the metadata and transcripts of the mixtures.
\begin{table}[ht]
\centering
\begin{tabularx}{\linewidth}{|@{}l >{\raggedright\arraybackslash}X >{\raggedright\arraybackslash}X@{}|}
\toprule
\textbf{Field} & \textbf{Description} & \textbf{Example} \\
\midrule
\texttt{clip\_id} & Unique identifier for each generated mixture. & \begin{tabular}[t]{@{}l@{}}\texttt{mix\_0000003\_0.00\_2\_}\\\texttt{7.4\_D}\end{tabular} \\
\hline
\texttt{audio\_path} & Relative path to the generated mixed waveform file. & \texttt{Output/audio/...wav} \\
\hline
\texttt{transcript} & Serialized list of speaker-attributed utterances with start and end timestamps. & \texttt{[(spk,text,t0,t1), ...]} \\
\hline
\texttt{overlap\_ratio\_target} & Target overlap ratio specified by the generation configuration. & \texttt{0.0} \\
\hline
\texttt{overlap\_ratio\_actual} & Overlap ratio measured from the generated mixture. & \texttt{0.0} \\
\hline
\texttt{speakers\_count} & Number of unique speakers. & \texttt{2} \\
\hline
\texttt{snr\_db} & Target signal-to-noise ratio in dB; can be empty for clean/no-noise settings. & \texttt{7.4}, \texttt{0.0}, \texttt{-5.0} \\
\hline
\texttt{noise\_type} & Noise category code used in sampling (e.g., Domestic/Public/Transport). & \texttt{D}, \texttt{P}, \texttt{T} \\
\hline
\texttt{overlap\_mask\_path} & Path to overlap mask data or placeholder if not generated. & \texttt{no mask} \\
\hline
\texttt{source\_files} & Serialized list of source speech files used to construct the mixture. & \begin{tabular}[t]{@{}l@{}}\texttt{['/Data/LibriSpeech/}\\\texttt{train-clean-100/8123/275209}\\\texttt{/8123-275209-0011.flac'}\\\texttt{, ...]} \end{tabular}\\
\hline
\texttt{noise\_files} & Serialized list of selected noise files (empty for clean conditions). & \texttt{'Data/Noise/TBUS/ch06.wav'} \\
\hline
\bottomrule
\end{tabularx}
\caption{Manifest schema used to store synthetic mixture metadata and transcripts.}
\label{tab:manifest_schema}
\end{table}


% ASR inference layer

% experimental setup for reproducibility



% metadata and transcription structure

% ASR inference layer

% experimental setup for reproducibility